% ECONOMICS_PROBLEM: {"id": "C79-271", "title": "Winning positions and nim-values in Circular Nim", "jel_code": "C79", "source_id": "OP-0220"}
% !TeX program = xelatex
\documentclass[11pt,a4paper]{article}
\usepackage[margin=23mm]{geometry}
\usepackage{fontspec}
\setmainfont{Latin Modern Roman}
\usepackage{amsmath,amssymb,mathtools}
\usepackage{xcolor,enumitem,booktabs,longtable,array,xurl,hyperref}
\definecolor{muted}{HTML}{53636C}
\hypersetup{colorlinks=true,urlcolor=blue,linkcolor=blue}
\setlength{\parindent}{0pt}
\setlength{\parskip}{5pt}
\newcommand{\R}{\mathbb R}
\newcommand{\E}{\mathbb E}
\newcommand{\N}{\mathbb N}
\newcommand{\ind}{\mathbf 1}
\newcommand{\Var}{\operatorname{Var}}
\newcommand{\Cov}{\operatorname{Cov}}
\newcommand{\supp}{\operatorname{supp}}
\DeclareMathOperator*{\argmin}{arg\,min}
\DeclareMathOperator*{\argmax}{arg\,max}
\newcommand{\entrymeta}[3]{{\sffamily\small\color{muted}\textbf{\texttt{#1}}\quad|\quad #2\par #3\par}}
\newcommand{\Problemref}[1]{\texttt{#1}}
\newcommand{\JELref}[2]{\texttt{#2} (JEL \texttt{#1})}
\begin{document}
\section*{Winning positions and nim-values in Circular Nim}
\textbf{Problem ID:} \texttt{C79-271}\quad \textbf{JEL:} \texttt{C79}\par
% BEGIN ECONOMICS STATEMENT
\label{problem:OP-0220}
% GAME_THEORY_PROBLEM: {"local_id": "GT-090", "original_number": 90, "kind": "P", "entry_kind": "statement_only", "published": false, "review_required": true, "category": "game_theory"}
\label{game:GT-090}
{\sffamily\small\color{muted}
\noindent\textbf{ID:} GT-090. \textbf{Category:} Impartial combinatorial games.
\textbf{Status:} P for the source's $CN(6,2)$ question; R for the wider classification program.
\par}
\subsection*{Definitions and mathematical statement}
Fix integers $n\geq2$ and $1\leq k\leq n$. A position is
$x=(x_0,\ldots,x_{n-1})\in\mathbb Z_{\geq0}^{n}$ with heap indices read modulo $n$. A legal move chooses $a\in\{0,\ldots,n-1\}$ and a vector $y$ satisfying
\[
 0\leq y_i\leq x_i\quad\text{for all }i,\qquad
 y_i=x_i\quad\text{if }i\notin\{a,a+1,\ldots,a+k-1\},\qquad y\neq x.
\]
In particular some selected heaps may remain unchanged. The game $CN(n,k)$ is played with alternating moves and normal play: a player with no move loses. It terminates because each move strictly decreases $\sum_i x_i$.

Let $\mathcal O_{n,k}(x)$ be the legal followers and define the Sprague--Grundy value by
\[
 g_{n,k}(x)=\operatorname{mex}\{g_{n,k}(y):y\in\mathcal O_{n,k}(x)\},
 \qquad \operatorname{mex}(A)=\min(\mathbb Z_{\geq0}\setminus A).
\]
The losing positions are $\mathcal P_{n,k}=\{x:g_{n,k}(x)=0\}$. The focused source question is to give an explicit structural characterization of
$\mathcal P_{6,2}$ and a winning move from each position outside it. The broader catalog program asks for such characterizations, and preferably formulas for $g_{n,k}$, in other unresolved parameter regimes.

\subsection*{Short English statement}
Classify losing positions, especially for six heaps with moves restricted to two adjacent heaps, and determine more of the game's nim-values.

\subsection*{Variants and scope boundaries}
The mex recurrence already gives an exhaustive finite algorithm for any fixed input. The research question concerns structural solution formulas, not decidability. Finding all zero nim-values is weaker than finding the entire Sprague--Grundy function. Mis\`ere play, requiring every selected heap to shrink, or deleting empty heap locations changes the game. For $k=1$ this is ordinary Nim; for $k=n$ only the all-zero position is losing.

\subsection*{Source relationship and status}
Section 5.4.2 of the inspected source singles out $CN(6,2)$ as an unresolved winning-strategy problem, referring to earlier Circular Nim work. It does not formulate one theorem asserting that every remaining $(n,k)$ case has a common closed form. The general classification and nim-value request is therefore an editorial expansion.

\subsection*{References}
\begingroup\footnotesize\raggedright
\noindent [S1] Ryuya Hora, \emph{Games as recursive coalgebras: A categorical view on the Nim-sum} (2025; inspected revision v3), arXiv:2510.22886v3. Inspected locator: Section 5.4.2, Circular Nim and the $CN(6,2)$ question.
\url{https://arxiv.org/html/2510.22886v3}
\par\endgroup

\subsection*{Common conventions: Source mathematical conventions}

\textbf{Finite games.} For a finite set $A$, $\Delta(A)$ is its probability
simplex. Players choose independent mixed strategies in Nash equilibrium;
correlation is allowed only when explicitly part of the solution concept.
In a game with payoff functions $u_i$ and strategy profile $x$, an additive
$\varepsilon$ Nash equilibrium satisfies
\[
 u_i(x)\geq u_i(x_i',x_{-i})-\varepsilon
 \quad\text{for every player $i$ and every admissible deviation $x_i'$}.
\]
For numerical approximation, payoffs are normalized as locally specified.
An $\varepsilon$ payoff residual is distinct from distance at most
$\varepsilon$ to the set of exact equilibria. Point convergence, convergence
of empirical play, and convergence in distance to an equilibrium set are
also distinct.

\textbf{Computational representations.} Unless stated otherwise, a complexity
question takes rational data with binary encoding; polynomial time is measured
in total bit length. A fixed-accuracy polynomial algorithm is not necessarily
polynomial in $1/\varepsilon$, and neither is necessarily polynomial in
$\log(1/\varepsilon)$. Succinct games and value, demand, or sampling oracles
must declare their interfaces. Exact real solutions require an explicit
representation; an unspecified list of arbitrary real numbers is not a
finite computational output. A claimed algorithmic barrier is meaningful
only for its specified model and reductions.

\textbf{Dynamic games.} Histories, information, observation, and allowable
randomization are part of the model. A stationary strategy depends only on
the current state; a positional strategy in a deterministic graph game
chooses one action at each controlled vertex. Nash equilibrium at an initial
state and equilibrium after every history are different requirements.
An expectation of a pathwise $\liminf$ payoff, a $\liminf$ of expected averages,
a limiting expected average, and a uniform finite-horizon equilibrium are
different criteria. None is silently substituted for another.

\textbf{Allocation and mechanisms.} A complete allocation assigns every item
exactly once unless otherwise stated. Zero-valued goods, negative values,
capacity constraints, feasibility, lotteries, and copies of goods affect
fairness definitions and are specified locally. Dominant-strategy incentive
compatibility, Bayesian incentive compatibility, universal truthfulness,
and truthfulness in expectation impose different inequalities. Whenever
an objective ratio has zero denominator, its convention is stated or those
instances are excluded.

\textbf{Infinite-dimensional models.} Expectations, integrals, stochastic
equations, and optimization objectives require their local regularity,
integrability, filtrations, and admissible domains. A backward--forward
mean-field system is a boundary-value problem; it is not an initial-value
semiflow without an additional construction. Long-time, large-population,
and vanishing-noise limits require an order of limits and a convergence
topology. If a minimum need not be attained, the objective is its infimum
or supremum and the approximation requirement is stated separately.
% END ECONOMICS STATEMENT
\end{document}
